% Part: sets-functions-relations
% Chapter: sets
% Section: pairs-and-products

\documentclass[../../../include/open-logic-section]{subfiles}

\begin{document}

\olfileid[fa-IR]{sfr}{set}{pai}
\olsection{زوج‌های مرتب، چندتایی‌ها و ضرب‌های دکارتی}

\begin{explain}
از اصل گسترش نتیجه می‌شود که اعضای یک مجموعه ترتیبی ندارند.
پس اگر بخواهیم ترتیب را بازنمایی کنیم، از \emph{زوج‌های مرتب}
$\tuple{x, y}$ استفاده می‌کنیم. در زوج نامرتب $\{x, y\}$، ترتیب
اهمیتی ندارد: $\{x, y\} = \{y, x\}$. اما در زوج مرتب اهمیت دارد:
اگر $x \neq y$، آنگاه $\tuple{x, y} \neq \tuple{y, x}$.

در نظریهٔ مجموعه‌ها چگونه باید به زوج‌های مرتب بیندیشیم؟ نکتهٔ اساسی
این است که می‌خواهیم این اندیشه را حفظ کنیم که دو زوج مرتب یکسان‌اند
اگر و تنها اگر مؤلفهٔ نخستشان یکسان و مؤلفهٔ دومشان نیز یکسان باشد؛ یعنی:
\[
  \tuple{a, b}= \tuple{c, d}\text{ اگر و تنها اگر هم }a = c
  \text{ و هم }b=d.
\]
می‌توانیم زوج‌های مرتب را در نظریهٔ مجموعه‌ها با تعریف
وینر--کوراتفسکی تعریف کنیم.
\end{explain}

\begin{defn}[زوج مرتب]\ollabel{wienerkuratowski}
	$\tuple{a, b} = \{\{a\}, \{a, b\}\}$.
\end{defn}

\begin{prob}
	با استفاده از \olref[sfr][set][pai]{wienerkuratowski} ثابت کنید
	$\tuple{a, b}= \tuple{c, d}$ اگر و تنها اگر هم $a = c$ و هم $b=d$.
\end{prob}

\begin{explain}
اکنون که تعریفی برای زوج مرتب تثبیت کرده‌ایم، می‌توانیم آن را برای
تعریف مجموعه‌های دیگری به کار بریم. برای نمونه، گاهی دنباله‌های مرتبِ
بیش از دو شیء را نیز می‌خواهیم، مانند \emph{سه‌تایی‌ها}
$\tuple{x, y, z}$، \emph{چهارتایی‌ها} $\tuple{x, y, z, u}$ و
همین‌طور ادامه. می‌توانیم سه‌تایی‌ها را زوج‌های مرتب خاصی بدانیم که
مؤلفهٔ نخستشان خود یک زوج مرتب است: $\tuple{x, y, z}$ همان
$\tuple{\tuple{x, y},z}$ است. دربارهٔ چهارتایی‌ها نیز چنین است:
$\tuple{x,y,z,u}$ همان $\tuple{\tuple{\tuple{x,y},z},u}$ است و
به همین ترتیب ادامه می‌یابد. به‌طور کلی، از \emph{$n$-تایی‌های مرتب}
$\tuple{x_1, \dots, x_n}$ سخن می‌گوییم.

برخی مجموعه‌های زوج‌های مرتب یا دیگر $n$-تایی‌های مرتب سودمند خواهند بود.
\end{explain}

\begin{defn}[ضرب دکارتی]
برای مجموعه‌های $A$ و $B$، \emph{ضرب دکارتیِ} آنها، $A \times B$،
چنین تعریف می‌شود:
\[
  A \times B = \Setabs{\tuple{x, y}}{x \in A \text{ و } y \in B}.
\]
\end{defn}

\begin{ex}
اگر $A = \{0, 1\}$ و $B = \{1, a, b\}$ باشند، ضرب آنها برابر است با
\[
A \times B = \{ \tuple{0, 1}, \tuple{0, a}, \tuple{0, b},
    \tuple{1, 1}, \tuple{1, a}, \tuple{1, b} \}.
\]
\end{ex}

\begin{ex}
اگر $A$ مجموعه‌ای باشد، ضرب $A$ در خودش، $A \times A$، را
$A^2$ نیز می‌نویسیم. این مجموعهٔ \emph{همهٔ} زوج‌های $\tuple{x, y}$
با $x, y \in A$ است. مجموعهٔ همهٔ سه‌تایی‌های $\tuple{x, y, z}$
برابر $A^3$ است و به همین ترتیب ادامه می‌یابد. می‌توانیم تعریفی
بازگشتی به دست دهیم:
\begin{align*}
  A^1 & = A\\
  A^{k+1} & = A^k \times A
\end{align*}
\end{ex}

\begin{prob}
همهٔ !!{element}s $\{1, 2, 3\}^3$ را فهرست کنید.
\end{prob}

\begin{prop}\ollabel{cardnmprod}
اگر $A$ دارای $n$ !!{element}s و $B$ دارای $m$ !!{element}s باشد،
آنگاه $A \times B$ دارای $n\cdot m$ عضو است.
\end{prop}

\begin{proof}
برای هر !!{element}~$x$ در~$A$، تعداد $m$ !!{element}s به صورت
$\tuple{x, y} \in A \times B$ وجود دارد. بگذارید
$B_x = \Setabs{\tuple{x, y}}{y \in B}$. چون هرگاه $x_1 \neq x_2$،
$\tuple{x_1, y} \neq \tuple{x_2, y}$، داریم
$B_{x_1} \cap B_{x_2} = \emptyset$. اما اگر
$A = \{x_1, \dots, x_n\}$، آنگاه
$A \times B = B_{x_1} \cup \dots \cup B_{x_n}$ و در نتیجه دارای
$n\cdot m$ !!{element}s است.

برای تجسم این مطلب، !!{element}s~$A \times B$ را در جدولی بچینید:
\[
\begin{array}{rcccc}
  B_{x_1} = & \{\tuple{x_1, y_1} & \tuple{x_1, y_2} & \dots & \tuple{x_1, y_m}\}\\
  B_{x_2} = & \{\tuple{x_2, y_1} & \tuple{x_2, y_2} & \dots & \tuple{x_2, y_m}\}\\
  \vdots & & \vdots\\
  B_{x_n} = & \{\tuple{x_n, y_1} & \tuple{x_n, y_2} & \dots & \tuple{x_n, y_m}\}
\end{array}
\]
ازآنجاکه همهٔ $x_i$ها با یکدیگر و همهٔ $y_j$ها با یکدیگر متفاوت‌اند،
هیچ دو زوجی در این جدول یکسان نیستند و شمار آنها $n\cdot m$ است.
\end{proof}

\begin{prob}
با استقرا روی~$k$ نشان دهید که برای هر $k \ge 1$، اگر $A$ دارای
$n$ !!{element}s باشد، آنگاه $A^k$ دارای $n^k$ !!{element}s است.
\end{prob}

\begin{ex}
اگر $A$ مجموعه‌ای باشد، \emph{واژه‌ای} بر روی~$A$ هر دنباله‌ای از
!!{element}s~$A$ است. می‌توان یک دنباله را $n$-تایی‌ای از
!!{element}s~$A$ دانست. برای نمونه، اگر $A = \{a, b, c\}$، آنگاه
دنبالهٔ ``$bac$'' را می‌توان سه‌تایی~$\tuple{b, a, c}$ دانست.
واژه‌ها، یعنی دنباله‌های نمادها، در علوم رایانه اهمیتی اساسی دارند.
طبق قرارداد، !!{element}s~$A$ را دنباله‌هایی به طول~$1$ و
$\emptyset$ را دنباله‌ای به طول~$0$ می‌شماریم. پس مجموعهٔ
\emph{همهٔ} واژه‌ها بر روی~$A$ برابر است با
\[
A^* = \{\emptyset\} \cup A \cup A^2 \cup A^3 \cup \dots
\]
\end{ex}

\end{document}
